\chapter{Kết quả thực nghiệm}
\label{chap:c4_results}

Dựa trên thiết kế ở Chương~3, chương này trình bày kết quả cài đặt, demo giao diện, kiểm thử và đánh giá hiệu quả chatbot RAG.

\section{Cài đặt các thành phần trọng tâm}
\label{sec:c4_key_impl}

Backend sử dụng DI tường minh qua factory functions trong \texttt{app.js}: mỗi module được khởi tạo bằng \texttt{buildXxxModule(\{models, services\})} với phụ thuộc truyền qua constructor. UnitOfWork pattern (\texttt{runInTransaction} + \texttt{lockRow} với \texttt{SELECT FOR UPDATE}) đảm bảo tính nguyên tử khi trừ tồn kho đồng thời. Hai cổng thanh toán VNPay (HMAC-SHA512) và MoMo (HMAC-SHA256) được tích hợp với cơ chế idempotency cho IPN callback. Hai cron job tự động dọn dẹp dữ liệu hàng ngày (2 giờ sáng) và hàng tuần (Chủ nhật 3 giờ sáng).

\section{Cài đặt chatbot AI: RAG Pipeline}
\label{sec:c4_chatbot_impl}

\subsection{Cài đặt UnifiedEmbeddingService}

\texttt{UnifiedEmbeddingService} triển khai chiến lược chain fallback qua ba nhà cung cấp embedding. Hàm \texttt{generateEmbedding(text, type)} kiểm tra danh sách providers ngay lúc khởi động -- nếu không có API key nào được cấu hình thì ném lỗi ngay lập tức thay vì để lỗi xuất hiện khi có request thực tế. Khi có providers, hàm duyệt tuần tự: provider nào thành công thì trả về vector ngay (nếu đang dùng fallback thì ghi log debug để theo dõi); provider nào thất bại thì ghi log cảnh báo và chuyển sang provider kế tiếp; provider cuối cùng thất bại thì ghi log lỗi và ném ngoại lệ để \texttt{hybridSearch} bắt và trả về mảng rỗng.

Mỗi provider xử lý phân biệt query\slash passage theo cơ chế riêng. Jina v3 nhận field \texttt{task} trong request body (\texttt{retrieval.query} hoặc \texttt{retrieval.passage}) và tự xử lý internally mà không cần tiền tố văn bản. HuggingFace e5-instruct yêu cầu tiền tố \texttt{passage:} cho đoạn văn và một instruction dài cho query (``\textit{Instruct: Given a product search query, retrieve relevant Vietnamese e-commerce products}''), còn e5-base (fallback cuối) dùng tiền tố \texttt{query:}/\texttt{passage:} đơn giản. Tất cả ba provider đều tạo ra vector 1024 chiều, đảm bảo tương thích với dữ liệu đã được index trước đó.

\subsection{Cài đặt HybridVectorStore}

Vector store được lưu dưới dạng file JSON (\texttt{data/vector-db.json}) thay vì database chuyên dụng như Pinecone hay Weaviate. Quyết định này phù hợp với quy mô dự án (dưới 10.000 sản phẩm) và tránh phụ thuộc vào dịch vụ ngoài, nhưng sẽ cần nâng cấp cho môi trường sản xuất quy mô lớn.

Hàm \texttt{hybridSearch} thực hiện tìm kiếm hai pha song song rồi gộp kết quả theo chiến lược kết hợp điểm từ hai phương pháp (score-based fusion). Kết quả từ hai nguồn được hợp nhất theo hai quy tắc:

\begin{equation}
\text{score}(p) =
\begin{cases}
\text{score\_dense}(p) + \Delta_{\text{overlap}} & \text{nếu } p \in \text{vector results} \\[4pt]
\text{minScore} + \dfrac{\text{score\_sparse}(p)}{\max\_kw} \times 0.05 & \text{nếu } p \text{ chỉ có trong keyword results}
\end{cases}
\label{eq:hybrid_score}
\end{equation}

Trong đó $\text{score\_dense}(p)$ là cosine similarity (chỉ tính nếu $\geq 0.45$), $\text{score\_sparse}(p)$ là điểm BM25-inspired tính bằng tổng trọng số token khớp (name $\times 3$, text $\times 1$) nhân với tỷ lệ token khớp trên tổng số token trong query, và $\Delta_{\text{overlap}} = 0.05$ nếu sản phẩm xuất hiện ở cả hai nguồn. Sản phẩm chỉ match keyword (không match semantic) được inject vào với score tối đa $0.45 + 0.05 = 0.50$ và được đánh dấu \texttt{lowConfidence=true}, luôn thấp hơn mọi kết quả vector tốt. Kết quả cuối cùng là tối đa 10 sản phẩm có điểm cao nhất (topK=10). Trường hợp không có sản phẩm nào vượt ngưỡng minScore=0.45, pipeline kích hoạt chế độ dự phòng: gọi lại \texttt{hybridSearch(finalQuery, 3, 0)} với minScore=0 để lấy tối đa 3 sản phẩm gần nhất và đánh dấu \texttt{lowConfidence=true}, cho phép LLM cảnh báo người dùng rằng kết quả có thể không chính xác.

\subsection{Cài đặt RAG Pipeline và ChatbotService}

Pipeline RAG 7 bước đã thiết kế ở Chương~3 (§\ref{sec:c3_rag_pipeline}) được cài đặt trong \texttt{chatbot-service.js}. Request đến \texttt{POST /api/chatbot/message} đi qua ba middleware (rate limiter, optional auth, Zod validation) trước khi vào pipeline. Các điểm cài đặt đáng chú ý so với thiết kế:

\begin{itemize}
  \item Bước 3: \texttt{isPromptInjection} cài đặt 24 regex (tuân thủ OWASP LLM01) bao gồm cả ký tự Unicode ẩn dùng để chèn lệnh.
  \item Bước 5: Cụm phủ định (``không cần iPhone'') được strip khỏi query trước khi embedding để tránh vector bị lệch về brand bị loại trừ.
  \item Bước 6: \texttt{parseLLMOutput} khớp sản phẩm qua 4 mức tin cậy giảm dần (exact name → version keyword → model number → 80\% từ), phát hiện hallucination khi LLM đề cập sản phẩm không có trong catalog.
  \item Khi LLM gặp lỗi tạm thời (429, 500, 503), hệ thống thử provider tiếp theo; lỗi cố định hoặc budget timer (30 giây) hết thì chuyển sang \texttt{simpleKeywordMatch}.
\end{itemize}

Session memory (\texttt{Map}, LRU 500 phiên, TTL 30 phút) và catalog cache (TTL 5 phút) được quản lý trong singleton \texttt{ChatbotService}, dọn dẹp tự động qua \texttt{\_evictStaleSessions()} sau mỗi lần cập nhật hội thoại.

\section{Kết quả demo giao diện}
\label{sec:c4_ui_demo}

Sau khi cài đặt hoàn tất, phần này trình bày kết quả giao diện thực tế của hệ thống. Frontend triển khai ba loại route guard (\texttt{ProtectedRoute}, \texttt{AdminRoute}, \texttt{PublicOnlyRoute}) và lazy loading 13 feature pages bằng \texttt{React.lazy()}. Floating chat widget hiển thị cố định trên mọi trang, cho phép mua hàng trực tiếp qua chat.

\subsection{Giao diện người dùng cuối}

Trang chi tiết sản phẩm (Hình~\ref{fig:screenshot_product_detail}) tập trung vào việc cung cấp đầy đủ thông tin để người dùng đưa ra quyết định mua hàng. Phần ảnh hiển thị ảnh sản phẩm lớn với bốn thumbnail gallery bên dưới để xem các góc độ khác nhau. Phần thông tin bên phải hiển thị tên sản phẩm (iPhone 17 Pro Max), giá gốc kèm badge giảm giá 5\%, nhãn tình trạng tồn kho, khu vực chọn biến thể dung lượng (256GB, 512GB, 1TB, 2TB) và màu sắc (Đen, Xanh đậm, Cam Vũ Trụ) theo thiết kế attribute system. Bộ chọn số lượng và nút ``Mua ngay'' được đặt nổi bật để giảm số thao tác đến khi đặt hàng. Phía dưới là hai tab ``Mô tả'' và ``Thông số kỹ thuật'' với bảng thông số chi tiết.

\begin{figure}[H]
  \centering
  \fitfig{figures/c4/screenshot_product_detail.png}
  \caption{Trang chi tiết sản phẩm với gallery ảnh, chọn biến thể và thông số kỹ thuật}
  \label{fig:screenshot_product_detail}
\end{figure}

\subsection{Giỏ hàng, thanh toán và đơn hàng}

Trang giỏ hàng (Hình~\ref{fig:screenshot_cart}) hiển thị danh sách sản phẩm với đầy đủ thông tin biến thể đã chọn (dung lượng, màu sắc), bộ điều chỉnh số lượng và nút xóa từng sản phẩm. Banner xanh phía trên thông báo ``Giỏ hàng đã đồng bộ với tài khoản'', xác nhận cơ chế merge cart đã hoàn thành sau khi đăng nhập. Panel ``Tóm tắt đơn hàng'' bên phải cho thấy mã giảm giá SUMMER2026 đã được áp dụng thành công với khoản giảm 21.995.000đ, hiển thị tạm tính, dòng giảm giá và tổng cộng rõ ràng. Hai nút hành động cuối trang là ``Tiến hành thanh toán'' và ``Tiếp tục mua sắm'' để quay lại cửa hàng.

\begin{figure}[H]
  \centering
  \fitfig{figures/c4/screenshot_cart.png}
  \caption{Trang giỏ hàng với mã giảm giá SUMMER2026 đã áp dụng và tóm tắt đơn hàng}
  \label{fig:screenshot_cart}
\end{figure}

Trang thanh toán (Hình~\ref{fig:screenshot_checkout}) trình bày form nhập thông tin giao hàng với các trường Tên, Họ, địa chỉ email, số điện thoại, và địa chỉ giao hàng ba cấp (Tỉnh/Thành phố, Quận/Huyện, Phường/Xã) thông qua dropdown phân cấp, cùng trường số nhà và tên đường. Địa chỉ đầy đủ được hiển thị tức thì ngay bên dưới form khi người dùng hoàn thành các trường địa chỉ. Panel tóm tắt đơn bên phải giữ nguyên thông tin mã giảm giá và tổng tiền từ trang giỏ hàng. Nút ``Tiếp tục thanh toán'' đưa người dùng đến bước tiếp theo để chọn phương thức thanh toán (COD, VNPay hoặc MoMo).

\begin{figure}[H]
  \centering
  \fitfig{figures/c4/screenshot_checkout.png}
  \caption{Trang thanh toán với form thông tin giao hàng địa chỉ ba cấp}
  \label{fig:screenshot_checkout}
\end{figure}

\subsection{Chatbot AI tư vấn sản phẩm}

Chat widget AI (Hình~\ref{fig:screenshot_chatbot}) xuất hiện cố định ở góc phải dưới màn hình trên mọi trang của ứng dụng. Khi người dùng nhấn vào bubble chat, cửa sổ ``Trợ lý mua sắm AI'' mở ra ngay trên trang hiện tại nhờ cơ chế React Portal, không làm gián đoạn việc duyệt web. Chatbot nhận câu hỏi tư vấn sản phẩm và trả về phản hồi kèm danh sách card gợi ý có thể cuộn ngang; mỗi card hiển thị ảnh thumbnail, tên sản phẩm, giá và nút thêm vào giỏ hàng trực tiếp từ chat mà không cần rời khỏi trang hiện tại.

\begin{figure}[H]
  \centering
  \fitfig{figures/c4/screenshot_chatbot.png}
  \caption{Chat widget AI tư vấn sản phẩm với card gợi ý kèm ảnh, giá và nút thêm giỏ hàng}
  \label{fig:screenshot_chatbot}
\end{figure}

Hình~\ref{fig:screenshot_chatbot_abbrev} minh họa khả năng xử lý viết tắt của chatbot. Hàm \texttt{expandAbbreviations} trong \texttt{ai-policy.js} sử dụng bảng \texttt{ABBREV\_MAP} gồm 16 mẫu regex cho viết tắt thương hiệu, model và hội thoại (trong tổng ABBREV\_MAP 71 pattern) để chuyển đổi viết tắt thông dụng trước khi đưa query vào vector search: các mẫu bao gồm ``ip'' (iPhone), ``ss'' (Samsung), ``mb'' (MacBook), ``pm'' (Pro Max), ``op'' (OPPO), ``rl'' (realme), ``r5'' (AMD Ryzen 5), ``r7'' (AMD Ryzen 7), ``bnh'' (bao nhiêu) và ``bh'' (bảo hành). Cơ chế này giúp pipeline RAG truy xuất đúng sản phẩm ngay cả khi người dùng nhắn theo phong cách viết tắt thông thường, ví dụ ``ip17 pm bnh'' được mở rộng thành ``iPhone 17 Pro Max bao nhiêu'' trước khi embedding.

\begin{figure}[H]
  \centering
  \fitfig{figures/c4/screenshot_chatbot_abbrev.png}
  \caption{Chatbot xử lý thành công câu hỏi viết tắt nhờ cơ chế expandAbbreviations với 71 mẫu regex}
  \label{fig:screenshot_chatbot_abbrev}
\end{figure}

Ngoài ra, chatbot từ chối lịch sự câu hỏi ngoài phạm vi (thời tiết, bóng đá, nấu ăn) nhờ hàm \texttt{isOffTopic} sử dụng regex thuần, hoàn thành trong dưới 10ms mà không tiêu tốn quota API.

\subsection{Giao diện quản trị}

Bảng điều khiển admin (Hình~\ref{fig:screenshot_admin_dashboard}) gồm hành động nhanh, card thống kê doanh thu và vận hành, biểu đồ doanh thu theo ngày với bộ lọc thời gian và xuất báo cáo.

\begin{figure}[H]
  \centering
  \fitfig{figures/c4/screenshot_admin_dashboard.png}
  \caption{Bảng điều khiển admin với hành động nhanh, thống kê tổng quan và biểu đồ doanh thu}
  \label{fig:screenshot_admin_dashboard}
\end{figure}

Các trang quản trị khác (sản phẩm, danh mục, thương hiệu, đơn hàng, người dùng, mã giảm giá, kho hàng) tuân theo cùng mẫu thiết kế: card thống kê phía trên, bảng danh sách với bộ lọc và phân trang, badge trạng thái trực quan. Điểm đáng chú ý là sự khác biệt giữa hai vai trò: khi quản trị viên đăng nhập, hệ thống hiển thị banner ``Chế độ xem'' màu xanh và ẩn các nút thao tác nghiệp vụ. Khi nhân viên bán hàng đăng nhập (Hình~\ref{fig:screenshot_staff_products}), banner biến mất và đầy đủ các thao tác CRUD được kích hoạt -- thể hiện trực quan mô hình RBAC đã thiết kế.

\begin{figure}[H]
  \centering
  \fitfig{figures/c4/screenshot_staff_products.png}
  \caption{Trang quản lý sản phẩm (nhân viên bán hàng) với nút thêm mới và đầy đủ 4 icon CRUD}
  \label{fig:screenshot_staff_products}
\end{figure}

\section{Kiểm thử hệ thống}
\label{sec:c4_testing}

\subsection{Chiến lược kiểm thử 5 tầng}

TechStore áp dụng chiến lược kiểm thử đa tầng theo mô hình Test Pyramid~\cite{fowler2012test_pyramid}, trong đó tầng thấp nhất có nhiều test nhất nhưng chạy nhanh nhất, tầng cao nhất có ít test hơn nhưng kiểm tra toàn bộ hệ thống. Hình~\ref{fig:testing_pyramid} minh họa kim tự tháp kiểm thử với tổng cộng ~7.303 test cases trên 325 suites, bổ sung bởi hai tầng chất lượng là mutation testing và property-based testing.

\begin{figure}[H]
  \centering
  \fitfig[0.90]{figures/c4/testing_pyramid.pdf}
  \caption{Kim tự tháp kiểm thử 5 tầng của hệ thống TechStore}
  \label{fig:testing_pyramid}
\end{figure}

Bảng~\ref{tab:test_results} trình bày kết quả kiểm thử chi tiết theo từng tầng.

\begin{table}[H]
  \centering
  \caption{Kết quả kiểm thử theo 5 tầng của hệ thống TechStore}
  \label{tab:test_results}
  \footnotesize
  \begin{tabular}{|l|p{3.5cm}|p{1.8cm}|p{1.3cm}|p{1.5cm}|p{2.2cm}|}
    \hline
    \textbf{Tầng} & \textbf{Framework} & \textbf{Số suites} & \textbf{Số test} & \textbf{Thời gian} & \textbf{Database} \\
    \hline
    BE E2E & Jest + Supertest & 5 & 100 & $\sim$22s & MySQL thật \\
    \hline
    BE API HTTP & Jest + Supertest & 39 & 675 & $\sim$160s & MySQL thật \\
    \hline
    BE Integration & Jest 29 & 38 & 210 & $\sim$57s & MySQL thật \\
    \hline
    BE Unit & Jest 29 & 215 & 5.381 & $\sim$12s & Mock (jest.fn) \\
    \hline
    FE Component & Jest + Testing Library & 28 & 937 & $\sim$14s & jsdom \\
    \hline
    \textbf{Tổng cộng} & & \textbf{325} & \textbf{$\sim$7.303} & & \\
    \hline
  \end{tabular}
  \normalsize
\end{table}

Tầng E2E (100 tests) kiểm tra 5 user flow hoàn chỉnh. API HTTP (675 tests) gửi real HTTP request qua Supertest. Integration (210 tests) kiểm tra với MySQL thật, bao gồm test race condition trừ tồn kho. Unit (5.381 tests) kiểm tra logic isolation với coverage 100\% statements/lines, 99,91\% branches. FE Component (937 tests) kiểm tra stores, utilities và React components. Test descriptions viết tiếng Việt, theo pattern Arrange-Act-Assert.

CI/CD qua GitHub Actions tự động lint, typecheck, build và unit test với coverage enforcement. Husky pre-commit hooks quét secret key và kiểm tra kiến trúc.

\section{Đánh giá hiệu quả chatbot RAG}
\label{sec:c4_rag_evaluation}

\subsection{Kịch bản đánh giá}

Để đánh giá hiệu quả của pipeline RAG, 20 kịch bản đại diện được thiết kế theo 5 nhóm yêu cầu khác nhau: tìm kiếm ngữ nghĩa (6 kịch bản), tên model cụ thể (4 kịch bản), viết tắt và sai chính tả (4 kịch bản), câu hỏi so sánh (3 kịch bản), và ngoài phạm vi (3 kịch bản). Bảng~\ref{tab:rag_scenarios} trình bày kết quả đánh giá.

\begin{table}[H]
  \centering
  \caption{Kết quả đánh giá chatbot RAG theo nhóm kịch bản}
  \label{tab:rag_scenarios}
  \footnotesize
  \begin{tabular}{|p{2.5cm}|c|p{3.5cm}|c|p{3.5cm}|}
    \hline
    \textbf{Nhóm} & \textbf{Số kịch bản} & \textbf{Ví dụ} & \textbf{Đúng} & \textbf{Ghi chú} \\
    \hline
    Tìm kiếm ngữ nghĩa & 6 & ``laptop học lập trình dưới 20 triệu'' & 5/6 & Dense retrieval tốt cho query mô tả \\
    \hline
    Tên model cụ thể & 4 & ``iphone 17 pro max bao nhiêu tiền'' & 4/4 & Hybrid Search bắt chính xác tên \\
    \hline
    Viết tắt\slash sai chính tả & 4 & ``ip17 pm 512gb gia bnh'' & 4/4 & \texttt{expandAbbreviations} xử lý đúng \\
    \hline
    Câu hỏi so sánh & 3 & ``so sanh macbook air vs pro'' & 2/3 & Đôi khi thiếu 1 sản phẩm gần ngưỡng \\
    \hline
    Ngoài phạm vi & 3 & ``thời tiết hôm nay thế nào'' & 3/3 & Từ chối đúng 100\%, $<$10ms \\
    \hline
    \textbf{Tổng} & \textbf{20} & & \textbf{18/20} & \textbf{Tỷ lệ chính xác 90\%} \\
    \hline
  \end{tabular}
  \normalsize
\end{table}

\textbf{Nhóm tìm kiếm ngữ nghĩa:} Khi người dùng mô tả nhu cầu bằng ngôn ngữ tự nhiên (``laptop chơi game cấu hình cao'', ``điện thoại pin trâu dưới 10 triệu''), dense retrieval với cosine similarity cho kết quả tốt vì vector embedding nắm bắt được ngữ nghĩa của câu mô tả. Chatbot truy xuất đúng loại sản phẩm trong 5 trên 6 trường hợp thử nghiệm trong nhóm này.

\textbf{Nhóm tên model cụ thể:} Khi người dùng gõ chính xác tên model (``iPhone 17 Pro Max 256GB'', ``Samsung Galaxy S24 Ultra''), sparse retrieval BM25 phát huy tác dụng vì khớp chính xác từ khóa tên thương hiệu và model. Overlap boost +0.05 giúp các sản phẩm xuất hiện ở cả hai nguồn được ưu tiên cao hơn.

\textbf{Nhóm viết tắt và sai chính tả:} \texttt{expandAbbreviations} xử lý thành công các viết tắt phổ biến như ``ip'' (iPhone), ``ss'' (Samsung), ``mb'' (MacBook), ``bnh'' (bao nhiêu), ``bh'' (bảo hành). Query rewriting bằng LLM xử lý các trường hợp phức tạp hơn như ``dtdd gia re pin trau'' (điện thoại di động giá rẻ pin trâu).

\textbf{Nhóm câu hỏi so sánh:} Đây là nhóm khó nhất vì chatbot cần truy xuất đủ thông tin cho cả hai sản phẩm. Kết quả đạt mức ``khá'' -- câu trả lời đúng nhưng đôi khi thiếu một trong hai sản phẩm nếu điểm similarity của một sản phẩm nằm sát ngưỡng 0.45.

\textbf{Nhóm ngoài phạm vi:} \texttt{isOffTopic} từ chối chính xác 100\% các câu hỏi về thời tiết, bóng đá, nấu ăn, sức khỏe, tin tức. Cơ chế từ chối sử dụng regex pattern thuần, không gọi LLM, nên phản hồi gần như tức thì (dưới 10ms).

\subsection{Phân tích hiệu năng pipeline}

Bảng~\ref{tab:rag_latency} trình bày thời gian xử lý trung bình của từng bước trong pipeline RAG dựa trên đo đạc thực tế.

\begin{table}[H]
  \centering
  \caption{Thời gian xử lý trung bình các bước trong RAG pipeline}
  \label{tab:rag_latency}
  \begin{tabular}{|l|p{5cm}|p{2cm}|p{3cm}|}
    \hline
    \textbf{Bước} & \textbf{Mô tả} & \textbf{Thời gian} & \textbf{Ghi chú} \\
    \hline
    validateMessage & Kiểm tra input hợp lệ & $<$1ms & CPU-bound, không I/O \\
    \hline
    expandAbbreviations & Mở rộng viết tắt & $<$1ms & Regex, không I/O \\
    \hline
    isOffTopic & Kiểm tra ngoài phạm vi & $<$10ms & Regex, không LLM \\
    \hline
    embed query (Jina) & Vector hóa câu hỏi & 300--600ms & API call đến Jina \\
    \hline
    hybridSearch & Tìm kiếm trong vector store & 30--80ms & File I/O + tính toán \\
    \hline
    rewriteQuery & LLM viết lại câu hỏi & 500--2000ms & API call, song song \\
    \hline
    augmentAndGenerate & LLM sinh phản hồi & 1000--4000ms & Phụ thuộc độ dài context \\
    \hline
    parseLLMOutput & Phân tích kết quả JSON & $<$5ms & CPU-bound \\
    \hline
    \textbf{Tổng (trung bình)} & & \textbf{2--5s} & Cảm nhận người dùng \\
    \hline
  \end{tabular}
\end{table}

Bước tốn thời gian nhất là sinh phản hồi từ LLM (1--4 giây), phụ thuộc vào tải của nhà cung cấp API và độ dài context. Nhờ chạy song song rewriteQuery và hybridSearch, thời gian xử lý tổng thể được rút ngắn đáng kể so với thực hiện tuần tự. Trong điều kiện bình thường, người dùng nhận được phản hồi sau 2--5 giây, cảm nhận này được giảm nhẹ bằng animation ``đang gõ'' (typing indicator) trong chat widget.

\subsection{Hạn chế của chatbot hiện tại}

Qua quá trình đánh giá, một số hạn chế được xác định. Thứ nhất, câu hỏi cần context nhiều lượt (``cái đó có màu trắng không?'' sau khi hỏi về iPhone) đôi khi không được xử lý đúng nếu session bị ngắt. Thứ hai, câu hỏi về chính sách bảo hành và đổi trả đôi khi không có thông tin đầy đủ trong vector store vì chỉ index thông tin sản phẩm, không index tài liệu chính sách. Thứ ba, với catalog sản phẩm nhỏ (dưới 100 sản phẩm trong môi trường demo), chatbot đôi khi gợi ý các sản phẩm không thực sự phù hợp do thiếu sự đa dạng.

\section{Mô tả đóng góp của sinh viên}
\label{sec:c4_contributions}

\subsection{Phần tự cài đặt}

Sinh viên tự nghiên cứu và triển khai: (1) toàn bộ pipeline RAG 7 bước với Hybrid Search, chain fallback embedding, session memory LRU -- đóng góp cốt lõi; (2) kiến trúc Modular Monolith 17 module với DI tường minh và EventBus; (3) RBAC 4 tác nhân thực thi 3 lớp (middleware, route guard, UI); (4) UnitOfWork + SELECT FOR UPDATE chống race condition; (5) tích hợp VNPay/MoMo với HMAC và idempotency; (6) bộ kiểm thử $\sim$7.303 test cases trên 5 tầng + mutation testing + property-based testing; (7) frontend Feature-Based 13 feature.

Dự án sử dụng các thư viện mã nguồn mở theo đúng mục đích thiết kế: Express + middleware ecosystem cho HTTP layer, Sequelize ORM cho database, shadcn/ui + Tailwind cho UI, TanStack Query + Zustand cho state management, và Jest + Supertest cho kiểm thử.

\bigskip
Chương này đã trình bày kết quả cài đặt, demo giao diện, kiểm thử 5 tầng và đánh giá hiệu quả chatbot RAG qua 20 kịch bản thực tế. Kết quả cho thấy pipeline RAG Hybrid xử lý tốt các truy vấn ngữ nghĩa, tên model cụ thể và viết tắt tiếng Việt, với thời gian phản hồi trung bình 2--5 giây. Phần Kết luận sẽ tổng kết kết quả đạt được và đề xuất hướng phát triển.
